\documentclass[10pt]{article}
\usepackage[textwidth=5.5in,textheight=9in,top=1in,left=1.5in,headheight=14pt]{geometry}
\usepackage[utf8]{inputenc}
\usepackage[T1]{fontenc}
\usepackage{times,microtype,amsmath,amssymb,graphicx,booktabs,natbib,xcolor,enumitem}
\usepackage[colorlinks=true,linkcolor=blue!60!black,citecolor=blue!60!black,urlcolor=blue!60!black]{hyperref}
\usepackage{titlesec}
\titlespacing*{\section}{0pt}{1.4ex plus 0.4ex}{0.8ex plus 0.2ex}
\titlespacing*{\subsection}{0pt}{1.0ex plus 0.3ex}{0.6ex plus 0.2ex}
\titleformat*{\section}{\fontsize{11}{13}\selectfont\bfseries}
\titleformat*{\subsection}{\fontsize{10}{12}\selectfont\bfseries}
\setlist{nosep, leftmargin=1.5em}
\usepackage{fancyhdr}
\pagestyle{fancy}
\fancyhf{}
\fancyhead[L]{\small D200 Machine Learning in Economics}
\fancyhead[R]{\small University of Cambridge}
\fancyfoot[C]{\thepage}
\renewcommand{\headrulewidth}{0.4pt}

\begin{document}
\thispagestyle{fancy}
\vspace{0.8em}
\begin{center}
{\fontsize{13}{16}\selectfont Does Welsh Media Need a Review?\\[4pt]Detecting Bias in Nation.Cymru's Political Reporting}\\[6pt]
{\small Cai Parry-Jones\\March 2026}
\end{center}
\vspace{0.8em}

\section{Introduction}

The run-up to the 2026 Welsh Parliament (Senedd) elections has been marked by growing accusations of political bias in Welsh media. Representatives from both the Conservatives and Reform UK have publicly questioned the impartiality of political news coverage in Wales.\footnote{BBC News, ``Reform calls for review into BBC and Plaid relationship'', 11 November 2025.}\textsuperscript{,}\footnote{Nation.Cymru, ``Tory MS raises concerns over `impartial' media coverage of Senedd election'', 21 February 2026.} But are these complaints supported by evidence?

I take a first step toward answering this by examining Nation.Cymru, a prominent Welsh political news outlet named Online News Site of the Year at the 2025 Wales Media Awards.\footnote{Journalists' Charity, Wales Media Awards 2025, 17 January 2026.} Its accessible API archive makes it a natural starting point for computational analysis. With polling placing Reform UK and Plaid Cymru as the two leading parties, I ask whether Nation.Cymru's coverage exhibits measurable differential sentiment bias between them. If bias is detectable in this outlet, it strengthens the case for a broader review of Welsh political media.

\section{Literature Review}

Recent systematic reviews confirm that media bias is both pervasive and increasingly detectable through modern natural language processing (NLP) methods \citep{rodrigogines2024}. For sentence-level bias detection, \citet{spinde2021} introduced the BABE expert-annotated dataset for sentence-level media bias detection, on which the \texttt{himel7/bias-detector} RoBERTa model used in this project was fine-tuned \citep{ghosh2025}. Recent work shows that LLMs can match or exceed fine-tuned classifiers on sentiment tasks, but at substantially higher computational cost \citep{shao2025}. Model cascading, i.e.,\ where a lightweight model handles simple tasks and a costlier model is reserved for harder cases, has been shown to reduce inference costs without a noticeable reduction in accuracy \citep{yue2024}. This motivates the two-stage design used for this project. No published computational studies of Welsh political media bias exist, making this the first such analysis.

\section{Data}

Nation.Cymru does not provide a structured database for research access. I therefore wrote a Python script to collect articles via its WordPress REST API, using \texttt{requests} for API calls and \texttt{BeautifulSoup} to parse the returned HTML content into plain text. To avoid requesting irrelevant articles as well as straining the Nation.Cymru server, I only request articles containing at least two Welsh politics keywords e.g.,\ `Senedd', political parties, notable political figures (see repository for full list). This yields 11{,}847 articles spanning 2017--2026.

I extract individual party mentions via string matching. For each mention, I extract a $\pm 2$ sentence context window and flag whether the text is quoted speech. The primary analysis restricts to Reform UK and Plaid Cymru, 2022-2026, news articles, non-quoted only, yielding 12{,}391 mentions (5{,}509 Reform; 6{,}882 Plaid). The secondary analysis covers four parties from 2022-2026 across news and opinion (49{,}661 mentions).

\section{Methodology}

\paragraph{Stage~1: Bias detection.}
Each mention's context window is classified as \textsc{biased}/\textsc{non-biased} using \texttt{himel7/bias-detector}, a RoBERTa model fine-tuned on the BABE dataset specifically for sentence-level media bias detection \citep{ghosh2025}. It is lightweight enough to classify all mentions locally, making it well suited as the first stage of the pipeline.

\paragraph{Stage~2: LLM sentiment attribution.}
A key challenge is distinguishing bias \textit{about} a party from sentiment \textit{expressed by} a party (e.g.,\ ``Plaid criticised the devastating policy'' has negative language, but Plaid is the speaker, not the target). I initially tested a dedicated sentiment model (\texttt{cardiffnlp/twitter-roberta-base-sentiment-latest}) but found it lacked the nuanced reasoning needed to distinguish speaker from target. I therefore use Claude Haiku 4.5 \citep{anthropic2025} via API as a prompted classifier, chosen as a cost-effective model capable of this task. For each biased mention it returns: (1)~\texttt{on\_target} (1/0): is the bias directed at the named party; (2)~sentiment score $s \in [-1, +1]$; (3)~one-sentence reasoning. The prompt includes speaker-vs-target examples and a self-check instruction. Off-target mentions are excluded.

\paragraph{Design choices.} Context window $\pm 2$ was selected after testing windows of 0, 1, and 2 sentences around each mention (\texttt{tune\_context\_window.py}); larger windows increased the bias detection rate, consistent with findings that transformer classifiers are sensitive to surrounding context for bias identification \citep{spinde2021}. Quoted speech is excluded to isolate editorial framing.

\paragraph{Implementation.} Stage~1: HuggingFace \texttt{transformers}, batch size 16, run locally on an M3 MacBook Air. Stage~2: Anthropic API, batches of 100; total API cost was approximately \$1.00 for the primary analysis and \$1.12 for the secondary (\$2.12 total, all other processing local), where LLM classification was capped at 250 biased mentions per party per article type. Without this cap, the secondary analysis alone would have cost an estimated \$10+. All code and data access instructions are available at \url{https://github.com/caitpj/d200_media_bias}.

\section{Results}

\subsection{Primary: Reform UK vs Plaid Cymru}

\begin{table}[ht]
\centering\small
\begin{tabular}{lcc}
\toprule
 & \textbf{Reform UK} & \textbf{Plaid Cymru} \\
\midrule
Mentions & 5{,}509 & 6{,}882 \\
On-target biased (\%) & 618 (11.2\%) & 226 (3.3\%) \\
Mean sentiment & $-0.721$ & $-0.242$ \\
\midrule
Bias rate $z$-test & \multicolumn{2}{c}{$z = 17.42$, $p < 0.001$, $h = 0.32$} \\
Sentiment $t$-test & \multicolumn{2}{c}{$t = -19.59$, $p < 0.001$, $d = -1.52$} \\
\bottomrule
\end{tabular}
\caption{Primary analysis: news coverage, 2022-2026.}
\label{tab:primary}
\end{table}

Reform receives on-target biased coverage at 3.4 times Plaid's rate (11.2\% vs 3.3\%), and 48.5\% of Reform's biased mentions are strongly negative ($s \leq -0.75$) versus 6.2\% for Plaid.

The LLM filtered 53.8\% of all biased mentions as off-target, with an asymmetric split: 63.5\% of Reform's biased mentions were on-target versus 26.5\% for Plaid. This reflects Plaid's frequent role as speaker rather than target. Manual review of a random sample of 50 LLM classifications found 1 error (2\% error rate).

\subsection{Secondary: all parties, news vs opinion}

\begin{figure}[ht]
\centering
\includegraphics[width=0.85\textwidth]{sentiment_bars.pdf}
\caption{Secondary analysis: mean on-target sentiment relative to article-type mean (news: $-0.527$; opinion: $-0.482$). Up to 250 biased mentions per party per article type were sent to the LLM; bars reflect only those classified as on-target (2022-2026). Error bars show 95\% confidence intervals.}
\label{fig:secondary}
\end{figure}

Figure~\ref{fig:secondary} shows sentiment relative to the article-type mean. Plaid Cymru receives notably more favourable coverage than the other parties across both news and opinion articles.

\section{Discussion}

\paragraph{Findings.} Nation.Cymru exhibits significant negative bias toward Reform UK compared to Plaid Cymru (Cohen's $d = -1.52$). However, this does not appear to be Reform-specific; the secondary analysis shows Plaid Cymru is the outlier, receiving more favourable coverage than the other parties. This result is consistent with the complaints discussed in the Introduction. On the methodology side, the LLM target-attribution step is important: without it, 53.8\% of biased mentions would be misattributed. This underscores the capacity of contemporary LLMs to handle nuanced classification tasks that exceed the capabilities of fine-tuned specialist models.

\paragraph{Limitations.} The analysis covers a single outlet, limiting generalisability. The bias detector was trained on US English news (the BABE dataset), not Welsh political text, so domain-specific language may be misclassified. LLM outputs are stochastic, so exact replication of results is not guaranteed. Keyword matching also misses references denoting previously mentioned parties (e.g.,\ ``the party is\ldots''). Statistical tests treat each mention as independent, but mentions from the same article may be correlated.

\paragraph{Next steps.} Extending to multiple outlets (BBC Wales, WalesOnline) would test whether the patterns found here reflect Nation.Cymru specifically or Welsh media more broadly. With a larger budget, a more capable LLM could improve classification accuracy, and the 250-mention sampling cap in the secondary analysis could be removed to cover all biased mentions. Taken further, the bias detection stage could be removed entirely, using the LLM for both bias detection and sentiment attribution in a single pass. Clustering standard errors by article would provide more robust inference.

\bibliographystyle{plainnat}
\begin{thebibliography}{6}
\bibitem[Anthropic(2025)]{anthropic2025} Anthropic. Claude {Haiku} 4.5 system card. Technical report, October 2025. \url{https://www.anthropic.com/claude-haiku-4-5-system-card}.
\bibitem[Ghosh et al.(2025)]{ghosh2025} H.~Ghosh, A.~Mosharafa, and G.~Groh. To bias or not to bias: Detecting bias in news with bias-detector. \textit{arXiv:2505.13010}, 2025.
\bibitem[Rodrigo-Gin\'{e}s et al.(2024)]{rodrigogines2024} F.-J.~Rodrigo-Gin\'{e}s, J.~Carrillo-de Albornoz, and L.~Plaza. A systematic review on media bias detection: What is media bias, how it is expressed, and how to detect it. \textit{Expert Systems with Applications}, 237:121641, 2024.
\bibitem[Shao(2025)]{shao2025} W.~Shao. Text sentiment classification optimization based on a fine-tuned {BERT} and large language model. \textit{Journal of Information Science}, 2025.
\bibitem[Spinde et al.(2021)]{spinde2021} T.~Spinde, M.~Plank, J.-D.~Krieger, T.~Ruas, B.~Gipp, and A.~Aizawa. Neural media bias detection using distant supervision with {BABE}. In \textit{Findings of EMNLP}, pages 1166--1177, 2021.
\bibitem[Yue et al.(2024)]{yue2024} M.~Yue, J.~Zhao, M.~Zhang, L.~Du, and Z.~Yao. Large language model cascades with mixture of thought representations for cost-efficient reasoning. In \textit{Proc.\ ICLR}, 2024.
\end{thebibliography}

\end{document}